\subsection{Orthonormal Coordinates and Inner-Product Identities}

Throughout this subsection, $\mathbb{F}$ is $\mathbb{R}$ or $\mathbb{C}$ and $\mathcal{V}$ is an inner-product space over $\mathbb{F}$.

\begin{theorem} \label{thm:orthonormal-basis-coordinates} \hlt{Coordinates in an Orthonormal Basis}\\
Let $\dim(\mathcal{V})=n\geq1$ and let $\beta=(\mathbf{e}_1,\ldots,\mathbf{e}_n)$ be an orthonormal basis. For $\mathbf{v}\in\mathcal{V}$:
\begin{enumerate}[label=(\roman*)]
\setlength{\itemsep}{0pt}
\item $\mathbf{v}=\sum\limits_{j=1}^{n}\innerprod{\mathbf{v},\mathbf{e}_j}\mathbf{e}_j$;
\item $\norm{\mathbf{v}}^2=\sum\limits_{j=1}^{n}\abs{\innerprod{\mathbf{v},\mathbf{e}_j}}^2$;
\item $[\mathbf{v}]_{\beta}=\bigl(\innerprod{\mathbf{v},\mathbf{e}_1},\ldots,\innerprod{\mathbf{v},\mathbf{e}_n}\bigr)^t$.
\end{enumerate}
\begin{proof}
\begin{enumerate}[label=(\roman*)]
\setlength{\itemsep}{0pt}
\item Write $\mathbf{v}=\sum\limits_{j=1}^{n}a_j\mathbf{e}_j$. For $k\in\{1,\ldots,n\}$, $\innerprod{\mathbf{v},\mathbf{e}_k}=\sum\limits_{j=1}^{n}a_j\innerprod{\mathbf{e}_j,\mathbf{e}_k}=a_k$ by orthonormality, which proves the expansion.
\item By (i) and Proposition \ref{prop:orthonormal-linear-combination-norm}, $\norm{\mathbf{v}}^2=\sum\limits_{j=1}^{n}\abs{\innerprod{\mathbf{v},\mathbf{e}_j}}^2$.
\item The coefficient computation in (i) means $[\mathbf{v}]_{\beta}=\bigl(\innerprod{\mathbf{v},\mathbf{e}_1},\ldots,\innerprod{\mathbf{v},\mathbf{e}_n}\bigr)^t$ by the coordinate-vector definition.
\end{enumerate}
\hlt{Structural Commentary}: Pairing with each basis vector extracts the coefficients before the norm and coordinate formulas follow.\\
\hlt{Key Insight}: Orthonormality turns coefficient extraction into an inner-product calculation.
\end{proof}
\end{theorem}

\begin{figure}[H]
  \centering
  \begin{tikzpicture}[
  >=Stealth,
  every node/.style={font=\small},
  axis/.style={gray!60, line width=0.4pt, ->},
  vec/.style={->, blue!75!black, line width=0.95pt},
  comp/.style={->, teal!70!black, line width=0.85pt},
  guide/.style={gray!45, dashed, line width=0.45pt},
  point/.style={circle, fill=black, inner sep=1pt}
]
  \coordinate (O) at (0,0);
  \coordinate (E1) at (3.2,0);
  \coordinate (E2) at (0,2.25);
  \coordinate (V) at (2.45,1.55);
  \coordinate (X) at (2.45,0);
  \coordinate (Y) at (0,1.55);

  \draw[axis] (-0.25,0) -- (3.65,0) node[right] {$\mathbf{e}_1$};
  \draw[axis] (0,-0.25) -- (0,2.65) node[above] {$\mathbf{e}_2$};
  \draw[guide] (X) -- (V) -- (Y);
  \draw[comp] (O) -- (X) node[midway, below] {$\innerprod{\mathbf{v},\mathbf{e}_1}\mathbf{e}_1$};
  \draw[comp] (X) -- (V) node[midway, right] {$\innerprod{\mathbf{v},\mathbf{e}_2}\mathbf{e}_2$};
  \draw[vec] (O) -- (V) node[midway, above left] {$\mathbf{v}$};

  \node[point, label=below left:{$\mathbf{0}$}] at (O) {};
  \node[point] at (V) {};
\end{tikzpicture}

  \caption{An orthonormal basis resolves a vector into perpendicular coordinate components.}
  \label{fig:orthonormal-coordinates}
\end{figure}

Figure \ref{fig:orthonormal-coordinates} illustrates Theorem \ref{thm:orthonormal-basis-coordinates}(i) for $n=2$: the two perpendicular coordinate components sum to $\mathbf{v}$.\par

\begin{remark} \label{rem:cross-terms} \hlt{Cross Terms}\\
The mixed terms $\innerprod{\mathbf{u},\mathbf{v}}$ and $\innerprod{\mathbf{v},\mathbf{u}}$ in $\norm{\mathbf{u}+\mathbf{v}}^2$ are called \hlt{cross terms}; $\norm{\mathbf{u}}^2$ and $\norm{\mathbf{v}}^2$ are the pure squared-length terms.
\end{remark}

\begin{theorem} \label{thm:pythagorean-inner-product} \hlt{Pythagorean Theorem}\\
$(\forall\,\mathbf{u},\mathbf{v}\in\mathcal{V})\bigl[\mathbf{u}\perp\mathbf{v}\Longrightarrow\norm{\mathbf{u}+\mathbf{v}}^2=\norm{\mathbf{u}}^2+\norm{\mathbf{v}}^2\bigr]$.
\begin{proof}
Fix $\mathbf{u},\mathbf{v}\in\mathcal{V}$ with $\mathbf{u}\perp\mathbf{v}$. Then $\innerprod{\mathbf{u},\mathbf{v}}=0$ by orthogonality and $\innerprod{\mathbf{v},\mathbf{u}}=0$ by Proposition \ref{prop:orthogonality-and-zero}(i).\\
Hence $\norm{\mathbf{u}+\mathbf{v}}^2=\innerprod{\mathbf{u}+\mathbf{v},\mathbf{u}+\mathbf{v}}=\norm{\mathbf{u}}^2+\innerprod{\mathbf{u},\mathbf{v}}+\innerprod{\mathbf{v},\mathbf{u}}+\norm{\mathbf{v}}^2=\norm{\mathbf{u}}^2+\norm{\mathbf{v}}^2$.\\
\hlt{Structural Commentary}: Orthogonality cancels the cross terms in the squared-norm expansion.\\
\hlt{Key Insight}: Perpendicular components add in squared length, not in length.
\end{proof}
\end{theorem}

\begin{proposition} \label{prop:pythagorean-converse-real-inner-products} \hlt{Pythagorean Converse over $\mathbb{R}$}\\
If $\mathbb{F}=\mathbb{R}$, then $(\forall\,\mathbf{u},\mathbf{v}\in\mathcal{V})\bigl[\mathbf{u}\perp\mathbf{v}\Longleftrightarrow\norm{\mathbf{u}+\mathbf{v}}^2=\norm{\mathbf{u}}^2+\norm{\mathbf{v}}^2\bigr]$.
\begin{proof}
Fix $\mathbf{u},\mathbf{v}\in\mathcal{V}$. Conjugate symmetry over $\mathbb{R}$ implies $\innerprod{\mathbf{v},\mathbf{u}}=\innerprod{\mathbf{u},\mathbf{v}}$, so $\norm{\mathbf{u}+\mathbf{v}}^2=\norm{\mathbf{u}}^2+2\innerprod{\mathbf{u},\mathbf{v}}+\norm{\mathbf{v}}^2$.
\begin{enumerate}[label=\roman*.]
\setlength{\itemsep}{0pt}
\item[$(\Rightarrow)$] If $\mathbf{u}\perp\mathbf{v}$, then $\innerprod{\mathbf{u},\mathbf{v}}=0$ by orthogonality, so the identity follows from the expansion.
\item[$(\Leftarrow)$] If the identity holds, then $2\innerprod{\mathbf{u},\mathbf{v}}=0$ by the expansion, so $\mathbf{u}\perp\mathbf{v}$ by orthogonality.
\end{enumerate}
\hlt{Structural Commentary}: The two directions compare the norm identity with its single real cross term.\\
\hlt{Key Insight}: Over $\mathbb{R}$, vanishing length excess is equivalent to orthogonality.
\end{proof}
\end{proposition}

\begin{theorem} \label{thm:bessel-inequality-finite-orthonormal-list} \hlt{Bessel's Inequality}\\
Let $m\geq1$ and let $(\mathbf{e}_1,\ldots,\mathbf{e}_m)$ be orthonormal in $\mathcal{V}$. Then $(\forall\,\mathbf{v}\in\mathcal{V})\,\sum\limits_{j=1}^{m}\abs{\innerprod{\mathbf{v},\mathbf{e}_j}}^2\leq\norm{\mathbf{v}}^2$.
\begin{proof}
Fix $\mathbf{v}\in\mathcal{V}$ and set $\mathbf{p}\defeq\sum\limits_{j=1}^{m}\innerprod{\mathbf{v},\mathbf{e}_j}\mathbf{e}_j$. For $k\in\{1,\ldots,m\}$, $\innerprod{\mathbf{v}-\mathbf{p},\mathbf{e}_k}=\innerprod{\mathbf{v},\mathbf{e}_k}-\sum\limits_{j=1}^{m}\innerprod{\mathbf{v},\mathbf{e}_j}\innerprod{\mathbf{e}_j,\mathbf{e}_k}=0$ by orthonormality.\\
Hence $\innerprod{\mathbf{v}-\mathbf{p},\mathbf{p}}=\sum\limits_{j=1}^{m}\overline{\innerprod{\mathbf{v},\mathbf{e}_j}}\innerprod{\mathbf{v}-\mathbf{p},\mathbf{e}_j}=0$, so $\mathbf{p}\perp(\mathbf{v}-\mathbf{p})$ by Proposition \ref{prop:orthogonality-and-zero}(i).\\
By Theorem \ref{thm:pythagorean-inner-product} and Proposition \ref{prop:orthonormal-linear-combination-norm}, respectively, $\norm{\mathbf{v}}^2=\norm{\mathbf{p}}^2+\norm{\mathbf{v}-\mathbf{p}}^2$ and $\norm{\mathbf{p}}^2=\sum\limits_{j=1}^{m}\abs{\innerprod{\mathbf{v},\mathbf{e}_j}}^2$. Thus $\sum\limits_{j=1}^{m}\abs{\innerprod{\mathbf{v},\mathbf{e}_j}}^2\leq\norm{\mathbf{v}}^2$.\\
\hlt{Structural Commentary}: The coordinate sum and its perpendicular residual split the squared norm by Pythagoras.\\
\hlt{Key Insight}: The residual norm is the nonnegative energy missing from selected coordinates.
\end{proof}
\end{theorem}

\begin{corollary} \label{cor:orthonormal-sequence-coefficients-vanish} \hlt{Coefficients Against an Orthonormal Sequence Vanish}\\
If $(\mathbf{e}_n)_{n\geq1}$ is orthonormal in $\mathcal{V}$, then $(\forall\,\mathbf{v}\in\mathcal{V})\,\lim\limits_{n\to\infty}\innerprod{\mathbf{v},\mathbf{e}_n}=0$.
\begin{proof}
Fix $\mathbf{v}\in\mathcal{V}$ and $\epsilon>0$, and set $S_{\epsilon}\defeq\{j\in\mathbb{N}_{\geq1}:\abs{\innerprod{\mathbf{v},\mathbf{e}_j}}\geq\epsilon\}$.\\
If $S_{\epsilon}$ were infinite, choose $N\in\mathbb{N}_{\geq1}$ with $N\epsilon^2>\norm{\mathbf{v}}^2$ and distinct $j_1,\ldots,j_N\in S_{\epsilon}$.\\
For this orthonormal sublist, $N\epsilon^2\leq\sum\limits_{k=1}^{N}\abs{\innerprod{\mathbf{v},\mathbf{e}_{j_k}}}^2\leq\norm{\mathbf{v}}^2$ by Theorem \ref{thm:bessel-inequality-finite-orthonormal-list}, contradicting $N\epsilon^2>\norm{\mathbf{v}}^2$. Thus $S_{\epsilon}$ is finite.\\
Choose $N_0\in\mathbb{N}_{\geq1}$ greater than every member of $S_{\epsilon}$, with $N_0=1$ when $S_{\epsilon}=\varnothing$. Then $(\forall\,n\geq N_0)\,\abs{\innerprod{\mathbf{v},\mathbf{e}_n}}<\epsilon$, so $\innerprod{\mathbf{v},\mathbf{e}_n}\to0$.\\
\hlt{Structural Commentary}: A contradiction argument applies Bessel to arbitrarily many allegedly large coefficients.\\
\hlt{Key Insight}: Finite total energy permits only finitely many coefficients above each threshold.
\end{proof}
\end{corollary}

\begin{theorem} \label{thm:parseval-identity-finite-dimensional} \hlt{Parseval's Identity}\\
Let $\dim(\mathcal{V})=n\geq1$ and let $\beta=(\mathbf{e}_1,\ldots,\mathbf{e}_n)$ be an orthonormal basis. For $\mathbf{v},\mathbf{w}\in\mathcal{V}$:
\begin{enumerate}[label=(\roman*)]
\setlength{\itemsep}{0pt}
\item $\innerprod{\mathbf{v},\mathbf{w}}=\sum\limits_{j=1}^{n}\innerprod{\mathbf{v},\mathbf{e}_j}\overline{\innerprod{\mathbf{w},\mathbf{e}_j}}$;
\item $\innerprod{\mathbf{v},\mathbf{w}}=\innerprod[\mathbb{F}^n]{[\mathbf{v}]_{\beta},[\mathbf{w}]_{\beta}}$.
\end{enumerate}
\begin{proof}
\begin{enumerate}[label=(\roman*)]
\setlength{\itemsep}{0pt}
\item Set $c_j\defeq\innerprod{\mathbf{v},\mathbf{e}_j}$ and $d_j\defeq\innerprod{\mathbf{w},\mathbf{e}_j}$ $(j\in\{1,\ldots,n\})$. Then $\mathbf{v}=\sum\limits_{j=1}^{n}c_j\mathbf{e}_j$ and $\mathbf{w}=\sum\limits_{k=1}^{n}d_k\mathbf{e}_k$ by Theorem \ref{thm:orthonormal-basis-coordinates}(i).\\
Hence $\innerprod{\mathbf{v},\mathbf{w}}=\sum\limits_{j=1}^{n}\sum\limits_{k=1}^{n}c_j\overline{d_k}\innerprod{\mathbf{e}_j,\mathbf{e}_k}=\sum\limits_{j=1}^{n}c_j\overline{d_j}$ by sesquilinearity and orthonormality.
\item By Theorem \ref{thm:orthonormal-basis-coordinates}(iii), $[\mathbf{v}]_{\beta}=(c_1,\ldots,c_n)^t$ and $[\mathbf{w}]_{\beta}=(d_1,\ldots,d_n)^t$. Therefore $\innerprod[\mathbb{F}^n]{[\mathbf{v}]_{\beta},[\mathbf{w}]_{\beta}}=\sum\limits_{j=1}^{n}c_j\overline{d_j}=\innerprod{\mathbf{v},\mathbf{w}}$ by (i).
\end{enumerate}
\hlt{Structural Commentary}: Basis expansions remove cross terms before the coordinate formula identifies the remaining sum.\\
\hlt{Key Insight}: Orthonormal coordinates preserve the inner product itself.
\end{proof}
\end{theorem}

\begin{theorem} \label{thm:matrix-entries-from-orthonormal-codomain} \hlt{Matrix Entries from an Orthonormal Codomain Basis}\\
Let $\dim(\mathcal{V})=n\geq1$ and $\dim(\mathcal{W})=m\geq1$, with $\mathcal{W}$ an inner-product space.\\
Fix a basis $\beta=(\mathbf{v}_1,\ldots,\mathbf{v}_n)$ of $\mathcal{V}$ and an orthonormal basis $\gamma=(\mathbf{w}_1,\ldots,\mathbf{w}_m)$ of $\mathcal{W}$. Then:
\begin{enumerate}[label=(\roman*)]
\setlength{\itemsep}{0pt}
\item $(\forall\,T\in\LL(\mathcal{V},\mathcal{W}),\ i\in\{1,\ldots,m\},\ j\in\{1,\ldots,n\})\,\bigl([T]_{\gamma\leftarrow\beta}\bigr)_{i,j}=\innerprod{T\mathbf{v}_j,\mathbf{w}_i}$;
\item $(\forall\,A\in\operatorname{Mat}(m,n,\mathbb{F}))(\exists!\,T\in\LL(\mathcal{V},\mathcal{W}))\,[T]_{\gamma\leftarrow\beta}=A$.
\end{enumerate}
\begin{proof}
\begin{enumerate}[label=(\roman*)]
\setlength{\itemsep}{0pt}
\item Fix $T\in\LL(\mathcal{V},\mathcal{W})$, $i\in\{1,\ldots,m\}$, and $j\in\{1,\ldots,n\}$. The $j$th column of $[T]_{\gamma\leftarrow\beta}$ is $[T\mathbf{v}_j]_{\gamma}=\bigl(\innerprod{T\mathbf{v}_j,\mathbf{w}_1},\ldots,\innerprod{T\mathbf{v}_j,\mathbf{w}_m}\bigr)^t$ by Theorem \ref{thm:orthonormal-basis-coordinates}(iii). Therefore its $i$th entry is $\innerprod{T\mathbf{v}_j,\mathbf{w}_i}$.
\item The map $\Phi:\LL(\mathcal{V},\mathcal{W})\to\operatorname{Mat}(m,n,\mathbb{F})$ defined by $\Phi(T)\defeq[T]_{\gamma\leftarrow\beta}$ is an isomorphism by Theorem \ref{thm:linear-maps-and-matrices-same-data}(i).\\
Hence every $A\in\operatorname{Mat}(m,n,\mathbb{F})$ has a unique preimage $T$ under $\Phi$.
\end{enumerate}
\hlt{Structural Commentary}: Coordinate extraction identifies entries; the earlier encoding isomorphism supplies existence and uniqueness.\\
\hlt{Key Insight}: Entry $(i,j)$ measures the $j$th basis image along output direction $i$.
\end{proof}
\end{theorem}

Figure \ref{fig:orthonormal-matrix-entries} visualizes Theorem \ref{thm:matrix-entries-from-orthonormal-codomain}(i): applying $T$ to $\mathbf{v}_j$ forms column $j$, and pairing with $\mathbf{w}_i$ reads entry $(i,j)$.

\begin{figure}[H]
  \centering
  \begin{tikzpicture}[
  x=0.88cm,
  y=0.88cm,
  >=Stealth,
  every node/.style={font=\small},
  space/.style={draw=gray!45, rounded corners=3pt, minimum width=3.25cm, minimum height=2.3cm},
  vec/.style={->, blue!75!black, line width=0.9pt},
  basis/.style={->, teal!70!black, line width=0.8pt},
  map/.style={->, gray!70, line width=0.7pt},
  measure/.style={->, orange!80!black, line width=0.65pt}
]
  \node[space] (V) at (0,0) {};
  \node[anchor=north west] at (-1.65,1.15) {$\mathcal{V}$};
  \node[space] (W) at (5.4,0) {};
  \node[anchor=north west] at (3.75,1.15) {$\mathcal{W}$};

  \coordinate (Ov) at (-0.95,-0.55);
  \coordinate (vj) at (0.85,0.65);
  \draw[vec] (Ov) -- (vj) node[midway, above left] {$\mathbf{v}_j$};
  \node[below] at (Ov) {$0$};

  \coordinate (Ow) at (4.55,-0.75);
  \coordinate (Tvj) at (6.45,0.75);
  \coordinate (wone) at (6.1,-0.75);
  \coordinate (wtwo) at (4.55,0.75);
  \draw[basis] (Ow) -- (wone) node[right] {$\mathbf{w}_1$};
  \draw[basis] (Ow) -- (wtwo) node[above] {$\mathbf{w}_2$};
  \draw[vec] (Ow) -- (Tvj) node[midway, above left] {$T\mathbf{v}_j$};
  \draw[gray!45, dashed] (Tvj) -- (6.45,-0.75);
  \draw[gray!45, dashed] (Tvj) -- (4.55,0.75);

  \draw[map] (1.95,0) -- (3.35,0) node[midway, above] {$T$};

  \node[draw=gray!45, rounded corners=3pt, align=center, inner sep=5pt] (coords) at (9.25,-0.05)
    {$[T\mathbf{v}_j]_{\gamma}$\\[2pt]
     $=\begin{bmatrix}
        \innerprod{T\mathbf{v}_j,\mathbf{w}_1}\\
        \vdots\\
        \innerprod{T\mathbf{v}_j,\mathbf{w}_m}
       \end{bmatrix}$};
  \draw[measure] (7.05,0.35) -- (coords.west);
\end{tikzpicture}

  \caption{The $j$th column of $[T]_{\gamma\leftarrow\beta}$ is obtained by resolving $T\mathbf{v}_j$ against the orthonormal basis of $\mathcal{W}$.}
  \label{fig:orthonormal-matrix-entries}
\end{figure}

\begin{corollary} \label{cor:inner-product-through-operator-matrix} \hlt{Inner Products Through an Operator Matrix}\\
Let $\dim(\mathcal{V})=n\geq1$ and let $\beta=(\mathbf{e}_1,\ldots,\mathbf{e}_n)$ be an orthonormal basis.\\
For $T\in\LL(\mathcal{V})$, set $A\defeq[T]_{\beta\leftarrow\beta}$. Then $(\forall\,\mathbf{v},\mathbf{w}\in\mathcal{V})\,\innerprod{T\mathbf{v},\mathbf{w}}=\innerprod[\mathbb{F}^n]{A[\mathbf{v}]_{\beta},[\mathbf{w}]_{\beta}}$.
\begin{proof}
Fix $\mathbf{v},\mathbf{w}\in\mathcal{V}$. The identity $[T\mathbf{v}]_{\beta}=A[\mathbf{v}]_{\beta}$ holds by Proposition \ref{prop:matrix-representation-action}. Hence $\innerprod{T\mathbf{v},\mathbf{w}}=\innerprod[\mathbb{F}^n]{A[\mathbf{v}]_{\beta},[\mathbf{w}]_{\beta}}$ by Theorem \ref{thm:parseval-identity-finite-dimensional}(ii).\\
\hlt{Structural Commentary}: Matrix action computes the first coordinate vector before Parseval transfers the inner product.\\
\hlt{Key Insight}: Orthonormal coordinates convert operator geometry into matrix algebra.
\end{proof}
\end{corollary}
